\section{Path and ordered-pair estimates}\label{q3:app:foundations}
This appendix proves the preliminary estimates used throughout. All functions and operators are on the space in Section~\ref{q3:sec:method}. Put $u_j=\E K^j\one$ and $G_t(z)=\sum_{j=0}^{t-1}z^j$, with $G_0=0$.

\subsection{Spectral bound and positive convolution}
\begin{lemma}\label{q3:lem:spectral}
We have
\[
 \norm K_{\op}\le\frac2\pi,\quad \norm K_{\op}^2\le\theta,\quad
 \norm{Kv}_\infty\le\norm v_2,\quad
 u_{2j}=(-1)^j\delta_{j},\quad u_{2j+1}=0,\quad \delta_{j+1}\le\theta \delta_{j}.
\]
Moreover $0\le d\le1/3$ and
\[
 \norm{G_t(K)}\le1\quad(t=0,1\text{ or }t\ge3),\qquad
 \norm{G_2(K)}^2\le1+\theta.
\]
\end{lemma}
\begin{proof}
For a real unit vector $v$ with $Kv\ne0$, set $w=Kv/\norm{Kv}_2$. Then $v,w$ are orthogonal unit vectors and $|\ip v{Kw}|=\norm{Kv}_2$. Write $v(x)+iw(x)=r(x)e^{i\varphi(x)}$, so $\E r^2=2$. Antisymmetrization gives
\[
 |\ip v{Kw}|\le\frac12\E_{x,y}r(x)r(y)|\sin(\varphi(y)-\varphi(x))|.
\]
Use the absolutely convergent Fourier series
\[
 |\sin t|=\frac2\pi-\frac4\pi\sum_{n\ge1}\frac{\cos(2nt)}{4n^2-1}.
\]
Each nonconstant average is the nonnegative square $|\E r e^{2in\varphi}|^2$ and has a negative coefficient. Hence the last upper bound is at most $(\E r)^2/\pi\le2/\pi$. The case $Kv=0$ is immediate. Also $4/\pi^2<13/32$; for instance Machin's identity and alternating arctangent truncation give $\pi>157/50$, and $320000<13\cdot24649$.

The row bound follows from Cauchy--Schwarz, and the moment identities from skew-adjointness. Expanding the nonnegative cyclic product on three independent points yields
\[
 0\le\E_{x,y,z}(1+K(x,y))(1+K(y,z))(1+K(z,x))=1-3d.
\]
The cubic expectation vanishes by a transposition. Thus $d\le1/3$.

For a rotation frequency $i\lambda$ put $x=\lambda^2\le\theta<1/2$. If $t$ is odd, $|G_t(i\lambda)|^2=(1+x^t)/(1+x)\le1$. If $4\mid t$, the numerator is $(1-x^{t/2})^2\le1$. If $t\equiv2\pmod4$ and $t\ge6$, use $2x^2+x^5\le17/32<1$ to obtain $(1+x^{t/2})^2\le1+x$. For $t=2$ the squared modulus is $1+x$. Orthogonal rotation blocks prove the assertions.
\end{proof}

Define $\beta_1=\beta_2=0$ and
\[
 \beta_n=\sum_{j=1}^{\lfloor(n-1)/2\rfloor}(-1)^{j-1}\delta_{j}\quad(n\ge3).
\]
The identity $S_n=AS_{n-1}+\beta_n\one$ proves \eqref{q3:eq:convolution} by iteration. Alternating sums give $d-\delta_{2}\le\beta_n\le d$, while $\E S_n=1-\beta_{n+1}$ for $n\ge1$. Consequently $0<f_n\le S_n$, $\mu_n\le1$, and $f_n,h_n\le2$. Also $\norm{S_n}_2\le1$ for $n\ge2$ by the geometric-sum bound, so $\E f_n^2\le1$ for $n\ge2$.

The initial means are
\[
 \mu_0=\mu_1=1,\quad \mu_2=1-d,\quad\mu_3=1-2d,\quad
 \mu_4=1-3d+\delta_{2},\quad\mu_5=1-4d+2\delta_{2}+d^2.
\]
They imply $\mu_n\le1-2d$ for $n=3,4,5$. For $n\ge6$, retain the four convolution terms with $f_0,f_1,f_2,f_3$ and use $\beta_j\ge d/2$. Then
\[
 \mu_n\le1-\frac d2-\frac d2(4-3d)\le1-2d.
\]
If $d=0$, $K\one=0$ and every rooted tree message equals one by leaf integration. Subsequent divisions by $d$ are therefore confined to $d>0$.

\subsection{Weighted propagation and a weighted path cube}
Put $N_\pm(F)=\E[(1\pm g)F^2]$. Row Cauchy--Schwarz gives $(AF)^2\le pA(F^2)$ and hence
\begin{equation}\label{q3:eq:energy}
 N_-(AF)\le\E[A^*(pq)F^2]\le N_-(F).
\end{equation}
Reversal gives $N_+(A^*F)\le N_+(F)$. The edge measure with density $1+K(x,y)$ has total mass one and marginals $p,q$. Cauchy--Schwarz and \eqref{q3:eq:energy} thus imply
\begin{equation}\label{q3:eq:bilinear}
 |\ip F{A^LG}|\le\sqrt{N_+(F)N_-(G)}\qquad(L\ge1).
\end{equation}
In particular $N_-(f_n)\le1$ and $\E(f_n^2/p)\le1$: the latter follows by dividing the row square inequality for $Af_{n-1}$ by $p$ and integrating.

\begin{lemma}\label{q3:lem:weightedpathcube}
For every $n\ge2$,
\[
 \E[qf_n^3]\le\E[qS_n^3]\le1-d/8.
\]
\end{lemma}
\begin{proof}
Write $x=g$, $y=S_n-1-g$. Skewness gives $\E y+\E(gy)=-d-u_{n+1}$ and
\begin{equation}\label{q3:eq:yvariance}
 \norm y_2^2-u_{n+1}\le \delta_{2}.
\end{equation}
For even $n$, this follows from the geometric-sum bound. For $n=2h+1$, integrate against the spectral measure of $g$. The integrand of the left side of \eqref{q3:eq:yvariance} minus the right side is
\[
 \frac{-t^2+(-1)^ht^h(1-t)+t^{2h+1}}{1+t}\le0.
\]
For odd $h$ the middle term is nonpositive and $t^{2h+1}\le t^2$. For even $h\ge2$, division by $t^2$ bounds the numerator by $-1+(1-t)+t=0$.

The pointwise identity
\[
 -2x^3-x^4-3x^2(1+x)y-3x^2y^2
 =x^2\left[1-\frac{(1+x)^2}{4}-3\left(y+\frac{1+x}{2}\right)^2\right]\le x^2
\]
then gives
\[
 \E[qS_n^3]\le1-2d+3(\norm y_2^2-u_{n+1})+\E[qy^3].
\]
For $n=2$ or $n\ge4$, $\norm y_\infty\le\sqrt d$ and $\norm y_2^2\le \delta_{2}$, so $\E[qy^3]\le(6/5)\delta_{2}$. The resulting bound is $1-(47/160)d\le1-d/8$. For $n=3$, $\norm y_2^2=\delta_{2}+\delta_{3}$, $u_4=\delta_{2}$, $\norm y_\infty^2\le3d/2$, and $\E[qy^3]\le9d/8$. Using $\delta_{3}\le d/4$ proves the remaining case.
\end{proof}

\subsection{Ordered pairs and exponential moments}
For $1\le a<b$, set $P_{ab}=h_af_b$, $R_{ab}=S_b-S_a$. Then
\begin{equation}\label{q3:eq:pairenvelope}
 P_{ab}\le e^{R_{ab}-\omega_{ab}},\qquad
 \omega_{ab}=
 \begin{cases}
 \beta_b,&a=1,2,\\
 2d-\delta_{2},&a=3,\\
 2d,&a\ge4.
 \end{cases}
\end{equation}
For $a=1,2$, use $f_a=S_a$, $h_a=2-S_a$ and $f_b\le S_b-\beta_b$. For $b>a\ge3$, subtraction of the positive convolutions gives
\[
 f_b-f_a=R_{ab}+\sum_{i=0}^{a-3}(\beta_{a-i}-\beta_{b-i})f_i
                 -\sum_{i=a-2}^{b-3}\beta_{b-i}f_i.
\]
The positive contribution in the middle sum is at most
$4\sum_{j\ge1}\delta_{2j}\le4\delta_{2}/(1-\theta^2)\le2d$; for $a=3$ it is at most $\delta_{2}$. Combine this with $h_a+f_b=2\mu_{a-1}+f_b-f_a$. Finally, for nonnegative $h,f$ with $h+f\le2+C$, AM--GM gives $hf\le(1+C/2)^2\le e^C$. This proves \eqref{q3:eq:pairenvelope}.

If a finite polynomial $Q$ has coefficients zero or one, grouping powers modulo four gives
\[
 |Q(i\lambda)|^2\le\frac{1+\lambda^2}{(1-\lambda^4)^2}
 \le\frac{32768}{16245}<C_*.
\]
For ordered disjoint pairs $a_1<b_1<\cdots<a_t<b_t$, put
\[
 P=\prod_{i}P_{a_ib_i},\quad R=\sum_iR_{a_ib_i},\quad
 \Delta=\sum_i\omega_{a_ib_i},\quad a=a_1.
\]
Then $R=K^aQ(K)g$ for a zero-one polynomial and
\begin{equation}\label{q3:eq:mask}
 P\le e^{R-\Delta},\qquad
 \norm R_\infty^2\le C_*\theta^{a-1}d,\qquad
 \norm R_2^2\le C_*\theta^ad.
\end{equation}
The same estimates are uniform in the number and size of the pairs. Also
\begin{equation}\label{q3:eq:weightedboundary}
 \E[pR_{ab}]=u_{a+1}-u_{b+1}.
\end{equation}
For disjoint pairs this implies $\E[pR]\le0$ if $a=1$, and $\E[pR]\le d$ if $a\ge2$. Indeed all boundary indices are distinct and $\sum_{j\ge2}\delta_{j}\le13d/19$. Under the original measure, $\E R\le d/2$ for $a\ge2$, and $\E R\le-d/2$ for $a=1$, because $\sum_{j\ge1}\delta_{2j}\le(416/855)d<d/2$.

For any probability measure and $|R|\le M$,
\begin{equation}\label{q3:eq:hoeffding}
 \E e^{sR}\le\exp\{s\E R+s^2M^2/2\}\quad(s\ge0).
\end{equation}
To see this, the second derivative of $\log\E e^{sR}$ is the tilted variance, at most $M^2$; integrate twice. Thus
\begin{equation}\label{q3:eq:criterion}
 \Delta-\E R\ge(s/2)M^2\quad\Longrightarrow\quad \E P^s\le1.
\end{equation}
One may use a scalar weighted probability measure in this criterion without changing any underlying operator or adjoint.

\begin{lemma}\label{q3:lem:pairmoments}
For every disjoint ordered pair list, $\E[pP^2]\le1$. If $a_1\ge3$, also $\E[pP^6]\le1$ and $\E P^6\le1$.
\end{lemma}
\begin{proof}
For one pair with $a\ge3$, \eqref{q3:eq:weightedboundary} gives $\E[pR_{ab}]\le \delta_{2}+\delta_{3}$. Thus $\Delta-\E[pR]\ge(1047/1024)d$, while $M^2\le(1+\theta)\theta^2d=(7605/32768)d$ and $3M^2\le(1047/1024)d$. For two or more pairs with $a_1\ge3$, $\Delta\ge(4-\theta)d$, $\E[pR]\le d$, and $M^2\le C_*\theta^2d$; these again satisfy \eqref{q3:eq:criterion} for $s=6$. The ordinary sixth moment follows by the same argument: for one pair $\E R\le \delta_{2}$ and for several pairs $\E R\le d/2$, which give stronger deficit bounds. H\"older gives the corresponding second moments.

For the weighted second moment with shortest length two and one pair, the following upper and lower bounds are valid, respectively:
\[
 \begin{array}{c|c|c}
 b&M^2\text{ upper bound}&\Delta-\E[pR]\text{ lower bound}\\\hline
 3&\delta_{2}&d+\delta_{2}\\4&\delta_{2}+\delta_{3}&d\\b\ge5&\delta_{2}&d-\delta_{2}-\delta_{3} .
 \end{array}
\]
The last comparison uses $1-2\theta-\theta^2=23/1024>0$. With at least two pairs, the deficit is at least $(2-\theta)d$ and $M^2\le C_*\theta d$.

For shortest length one and one pair, $b=2,3,b\ge4$ give $(M^2,\Delta-\E[pR])$ bounded by $(d,d)$, $(d+\delta_{2},2d+\delta_{2})$ and $(d,2d-\delta_{2}-\delta_{3})$. If there are exactly two pairs and the first is $(1,2)$, the deficit is at least $(2071/1024)d>C_*d\ge M^2$. For at least three pairs beginning with $(1,2)$ it is at least $(4-\theta)d$. For at least two pairs with first long arm at least three it is at least $(3-\theta)d$. These exhaust the cases and satisfy \eqref{q3:eq:criterion} with $s=2$ under $p\mu$.
\end{proof}

\subsection{Short-pair cubic estimates}\label{q3:sec:shortpaircubes}
\begin{lemma}\label{q3:lem:shortpaircubes}
For $n\ge2$ and $m\ge3$,
\[
 \E(qf_n)^3\le1-d/6,\qquad \E(h_2f_m)^3\le1-d/10.
\]
In particular,
\begin{equation}\label{q3:eq:quantshort}
 \E(qf_2)^3\le1-\alpha_*d,\qquad
 \alpha_*=\frac{66751}{92160}>\frac{18}{25}.
\end{equation}
\end{lemma}
\begin{proof}
If nonnegative $u,v$ satisfy $u+v\le2m+R$, then $(uv)^3\le(m+R/2)^6$. For $m\ge0$, $|R|\le M$, expansion gives
\[
 \E(m+R/2)^6\le m^6+3m^5\E R+\mathcal C(m,M)\E R^2,
\]
\[
 \mathcal C(m,M)=\frac{15}4m^4+\frac52m^3M+
 \frac{15}{16}m^2M^2+\frac3{16}mM^3+\frac1{64}M^4.
\]
We will use $\mathcal C\le16133/2880$ for $m\le1,M\le3/5,M^2\le1/3$; $\mathcal C\le1579/256$ for $m\le1,M\le3/4,M^2\le1/2$; and $\mathcal C<197/40$ for $m\le15/16$ under the latter two bounds on $M$. These follow term by term from the displayed polynomial.

For $qf_2$, use $q+f_2=2+Kg$, $\E Kg=-d$, $\E(Kg)^2=\delta_{2}$, $M\le\sqrt d$. The bound is
$1-3d+(16133/2880)\delta_{2}\le1-\alpha_*d$.
For $qf_3$, take $m=1-d/2$, $R=Kg+K^2g$, so $\E R=-d$, $\E R^2=\delta_{2}+\delta_{3}\le3d/5$, $M\le3/4$. The identity
\[
 (1-t/2)^6-3t(1-t/2)^5
 =1-6t+\frac{45}4t^2-10t^3+\frac{75}{16}t^4-\frac98t^5+\frac7{64}t^6
\]
is at most $1-6t+45t^2/4-9t^3$ on $[0,1/5]$ and at most the analogous expression with $-8t^3$ on $[0,1/3]$. Splitting at $d=1/6$ gives deficits $3d/8$ and $331d/1800$, both at least $d/6$.
For $n\ge4$, set $\tau=(19/32)d$, $m=1-\tau/2$ and $R=S_n-1-g$. Positive convolution gives $q+f_n\le2m+R$, while $\E R\le-\tau$, $\E R^2\le \delta_{2}$, $M\le\sqrt d$. The preceding polynomial bounds yield a deficit coefficient at least
\[
 \frac{80275}{32768}-\frac{209729}{92160}
 =\frac{256711}{1474560}>\frac16.
\]

For $h_2f_3$, use $h_2+f_3=2-d+K^2g$, $M\le3/8$, $\E R=0$, $\E R^2=\delta_{3}$. Since $\mathcal C\le5$ and $(1-d/2)^6\le1-7d/4$, the deficit is $(947/1024)d$.
For $h_2f_4$, use
\[
 m=1-d,\quad R=K^2g+K^3g-dg,\quad
 \E R=\delta_{2},\quad \E R^2=\delta_{3}+\delta_{4}+2d\delta_{2}+d^3.
\]
When $d\le1/6$, $M\le1/2$, $\E R^2<2d/5$, $\mathcal C<16/3$, and the constant and mean terms are at most $1-73d/32$. The deficit is at least $71d/480$. When $d\ge1/6$, $M<4/5$, $\mathcal C<7/2$, $\E R^2<5d/8$. The constant and mean terms are bounded by $(1-d)^5(1+7d/32)\le1-5d/2$: this function is convex and its value at $1/3$ is $103/729<1/6$. The deficit is then at least $5d/16$.

For $h_2f_n$, $n\ge5$, positive convolution gives $h_2+f_n\le2-2\tau+R$ with
\[
 R=S_n-S_2-\tau g,\quad
 \E R\le \delta_{2},\quad \E R^2\le(\theta+(19/32)d)^2d,\quad
 M\le23/40,\quad M^2<1/3.
\]
For $d\le2/19$, the constant and mean terms are at most $1-915d/512$, and the variance term at most $80665d/65536$; their difference is $36455d/65536$. For $d\ge2/19$, $m=1-\tau\le15/16$, $\mathcal C<9/2$, and the variance term is at most $841d/512$. Convexity of
$W(d)=(1-(19/32)d)^5(1+5d/8)$ and $W(1/3)<5/12$ give $W(d)\le1-7d/4$. The remaining deficit is $55d/512>d/10$. All interval endpoints are included, and $d=0$ gives equality.
\end{proof}

\subsection{Every tree with at most two branch vertices}\label{q3:app:spiderboundary}
Every rooted spider has an orientation with $N_-(F)\le1$. To prove this, first restore equal pairs only after estimating the residual distinct list $\ell_1<\cdots<\ell_t$. For $t=0$ or $1$, use $F=1$ or $f_{\ell_1}$ and \eqref{q3:eq:energy}. For positive even $t$, use the reversal of the canonical pair product and Lemma~\ref{q3:lem:pairmoments}. For odd $t\ge3$ with $\ell_1=1$, put $F=pP$, with the remaining pairs positive canonical. Then
\[
 N_-(F)=\E[p(1-g^2)P^2]\le\E[pP^2]\le1.
\]
For odd $t\ge3$ with $\ell_1\ge2$, put $F=f_{\ell_1}P(-K)$, where $P$ pairs the remaining arms. Their shortest length is at least three, so
\[
 N_-(F)\le(\E[qf_{\ell_1}^3])^{2/3}
          (\E[qP(-K)^6])^{1/3}\le1.
\]
Use Lemmas~\ref{q3:lem:weightedpathcube} and \ref{q3:lem:pairmoments}. Complete reversal gives an $N_+$ certificate.

For two branch vertices choose the left endpoint with $N_+(F)\le1$, the right with $N_-(G)\le1$, and direct the entire connector left to right. Exact integration and \eqref{q3:eq:bilinear} give $Z_D=\ip F{A^LG}\le1$. For one branch vertex, designate one actual arm as an inward path and use the remaining root certificate in the same inequality. A path uses $\mu_n\le1$ and a single vertex has density one. This proves the lower-branch cases used in Section~\ref{q3:sec:assembly}.
